

\documentclass[letter, 11pt]{article}

    %%%%%%%%%%%%%%%%%%%%%%%%%%%%%%
    \oddsidemargin  0.0in
    \evensidemargin 0.0in
    \textwidth      6.8in
    \headheight     0.0in
    \topmargin      0.0in
    \textheight=9.0in
    \parindent=0in
    %%%%%%%%%%%%%%%%%%%%%%%%%%%%%%

\renewcommand{\familydefault}{\sfdefault} % Police 

\usepackage[dvipsnames]{xcolor}
\usepackage{amsthm,amsfonts,amssymb,amsmath, enumitem}
\usepackage{graphicx}   %To include graphics
\usepackage[small,compact]{titlesec}

\usepackage{textcomp}
\newcommand{\mytilde}{$\sim$}
\usepackage{fancyvrb}
\fvset{framesep=1mm,fontfamily=courier,framerule=.3mm,numbersep=1mm,commandchars=\\\{\}}

\usepackage[utf8]{inputenc}

\usepackage[french]{babel}

\usepackage{fancyhdr}
\pagestyle{fancy}

\usepackage{version} % Pour faire un texte troué pour les étudiants
%\includeversion{etudiant}\excludeversion{professeur} %version etudiant
\includeversion{professeur}\excludeversion{etudiants} %version prof
\newcommand{\etudiant}[1]{\begin{etudiants}#1 \end{etudiants}}
\newcommand{\prof}[1]{\begin{professeur}\color{blue} {\bf Solution :} #1 \end{professeur}}


\fancyhead[L]{STT-2100 \\ Régression}
\fancyhead[R]{A.-S. Charest, E. Reny-Nolin, A.-A. Vallée \\ Université Laval}


\begin{document}

\vspace*{0.01 cm}
\begin{center}
{\Large Ajuster une RLM avec R }
\end{center}

Les mots en rouge indiquent les noms choisis par l'utilisateur.\\


On utilise la fonction {\tt lm( )} (pour {\it linear model}), comme en RLS :

\begin{Verbatim}[commandchars=\\\{\}]
{\textcolor{red}{mod}} <- lm({\textcolor{red}{varY}} \mytilde {\textcolor{red}{varX1}} + {\textcolor{red}{varX2}} + {\textcolor{red}{varX3}}, data = {\textcolor{red}{donnees}})
\end{Verbatim}


Cette fonction retourne un objet de type liste et de classe lm, contenant plusieurs composantes dont les suivantes :
\begin{itemize}\itemsep-1mm
\item \texttt{{\textcolor{red}{mod}}\$coefficients} : un vecteur contenant les coefficients ;
\item \texttt{{\textcolor{red}{mod}}\$residuals} : un vecteur contenant les résidus ;
\item \texttt{{\textcolor{red}{mod}}\$fitted.values} : un vecteur contenant les valeurs prédites.
\end{itemize}

Les fonctions {\tt summary({\textcolor{red}{mod}})} et {\tt anova({\textcolor{red}{mod}})} produisent des informations importantes, dont les tests $t$ partiels sur les coefficients et certains tests $F$.

\vspace{0.1 in}
\subsection*{Retrouver les résultats par calcul matriciel}

Supposons qu'on dispose d'un tableau de données contenant les 20 observations de $(x_1,\,x_2,\,x_3,\,Y)$.

\begin{Verbatim}
{\textcolor{ForestGreen}{# Création de la matrice X et du vecteur Y}}
{\textcolor{red}{X}} <- as.matrix(cbind(rep(1,20),donnees[c("{\textcolor{red}{varX1}}","{\textcolor{red}{varX2}}","{\textcolor{red}{varX3}}")]))
{\textcolor{red}{Y}} <- as.matrix(donnees["{\textcolor{red}{varY}}"])
{\textcolor{red}{n}} <- dim(X)[1]
{\textcolor{red}{pprime}} <- dim(X)[2]

{\textcolor{ForestGreen}{# Création de la matrice H}}
{\textcolor{red}{H}} <- {\textcolor{red}{X}} %*% solve( t({\textcolor{red}{X}}) %*% {\textcolor{red}{X}} ) %*% t({\textcolor{red}{X}})

{\textcolor{red}{Yhat}} <- {\textcolor{red}{H}} %*% {\textcolor{red}{Y}}  {\textcolor{ForestGreen}{# Calcul des valeurs prédites}}
{\textcolor{red}{ehat}} <- {\textcolor{red}{Y}} - {\textcolor{red}{Yhat}} {\textcolor{ForestGreen}{# Calcul des résidus}}

{\textcolor{red}{sigma2hat}} <- t({\textcolor{red}{ehat}}) %*% {\textcolor{red}{ehat}} / ( {\textcolor{red}{n}} - {\textcolor{red}{pprime}})
\end{Verbatim}





\end{document}

